% Generated draft: requires mathematical and source review.
\documentclass[11pt]{article}
\usepackage[margin=25mm]{geometry}
\usepackage{fontspec}
\setmainfont{Latin Modern Roman}
\usepackage{amsmath,amssymb,mathtools,unicode-math}
\setmathfont{Latin Modern Math}
\usepackage{enumitem,needspace,xurl,hyperref}
\setlength{\parindent}{0pt}
\setlength{\parskip}{5pt}
\setlength{\emergencystretch}{3em}
\newcommand{\N}{\mathbb{N}}
\newcommand{\Z}{\mathbb{Z}}
\newcommand{\Q}{\mathbb{Q}}
\newcommand{\R}{\mathbb{R}}
\newcommand{\C}{\mathbb{C}}
\newcommand{\F}{\mathbb{F}}
\newcommand{\A}{\mathbb{A}}
\newcommand{\PP}{\mathbb{P}}
\newcommand{\E}{\mathbb{E}}
\newcommand{\eps}{\varepsilon}
\newcommand{\dd}{\,\mathrm{d}}
\begin{document}
\section*{UnsolvedMath ID: 30006622}
\textbf{Source problem code:} catalogue entry
\subsection{Definitions and mathematical statement}
\textbf{Definitions and conventions.} A hyperbolic group means a finitely generated group whose Cayley graph is Gromov hyperbolic. A subgroup is quasiconvex if geodesics between its elements remain within a bounded distance of the subgroup in a Cayley graph. A cubical action is understood below as an action by cubical automorphisms on a CAT(0) cube complex, equipped with its standard piecewise Euclidean metric. Such an action is cobounded if some orbit has a bounded neighbourhood equal to the whole complex. For a hyperplane $H$, its stabilizer is $\operatorname{Stab}_G(H)=\{g\in G:gH=H\}$. These conventions make the extracted questions readable, but the supplied original excerpt does not specify the intended cube complexes or metrics.

A group is cocompactly cubulable if it admits a proper, cocompact cubical action on a CAT(0) cube complex. It is a-T-menable if it admits a metrically proper affine isometric action on a real Hilbert space. It is residually finite if, for every $g\ne 1$, there is a homomorphism $\phi:G\to F$ to a finite group with $\phi(g)\ne 1$. In the unrestricted sense, it is linear if there exist a field $K$, an integer $n\geq 1$, and an injective homomorphism $G\to\operatorname{GL}_n(K)$. The record does not specify a field or characteristic for linearity.

For a metric space $(M,d)$, define
\[
 I_M(x,y)=\{z\in M:d(x,y)=d(x,z)+d(z,y)\}.
\]
The space is median if, for every $x,y,z\in M$, the intersection
\[
 I_M(x,y)\cap I_M(y,z)\cap I_M(z,x)
\]
contains exactly one point, denoted $m(x,y,z)$. Its median rank is finite if there is a finite bound on the integers $n$ for which the coordinatewise median algebra $\{0,1\}^n$ admits an injective median-preserving map into $M$.

For an isometric action of a discrete group, metric properness means that $\{g\in G:d(x,gx)\leq R\}$ is finite for every $x$ and every finite $R\geq 0$. Cocompactness means that the orbit space is compact. The supplied record uses the word ``proper'' without identifying whether metric properness or a topological properness condition is intended; those conditions need not agree for arbitrary median spaces.

\textbf{Questions extracted in the curated statement.} The following are separate questions, rather than a single assertion with a disjunctive conclusion.
\begin{itemize}
\item[(A)] Let $G$ be a hyperbolic group. Suppose that for every infinite-order element $g\in G$ there exists a cobounded cubical action $G\curvearrowright X_g$ such that every hyperplane stabilizer is quasiconvex in $G$ and
\[
 \sup_{n\in\Z}d_{X_g}(x,g^n x)=\infty
\]
for some, equivalently every, $x\in X_g$. Must $G$ be cocompactly cubulable? Must $G$, at least, be a-T-menable? Here the action is allowed to depend on $g$, following the natural reading of ``every infinite-order element \ldots\ in some \ldots\ action.'' The original excerpt does not establish this quantifier order or supply the omitted hypotheses.
\item[(B)] Let $G$ be a hyperbolic group admitting a proper, cocompact isometric action on a finite-rank median space $M$. Must $G$ be residually finite? Must $G$ be linear? The record imposes no explicit completeness or local compactness assumption on $M$, and does not specify the intended meaning of properness or the field for linearity.
\item[(C)] Let $(X,d)$ be a complete, simply connected path-metric space, where path-metric means that distance is the infimum of lengths of joining curves. Suppose every point has a neighbourhood in which every triple has a unique median, without assuming a uniform neighbourhood radius. Must $(X,d)$ be median?

One precise possible reading of the local hypothesis is
\[
 \forall p\in X\ \exists U_p\subseteq X\text{ open},\quad p\in U_p,\quad (U_p,d|_{U_p\times U_p})\text{ is median}.
\]
Equivalently, under this reading, every triple in $U_p$ has exactly one point in $U_p$ satisfying the three metric interval conditions for the restricted metric. The source does not settle whether neighbourhoods instead carry their intrinsic path metrics, whether medians are required to be unique in the entire ambient space, or what precise uniform local hypothesis is being weakened. Consequently this displayed hypothesis is an explicitly qualified interpretation. ``Without a uniform radius'' means that no common positive radius is assumed; it does not require the failure of every uniform-radius condition.
\end{itemize}

\textbf{Unresolved scope of the supplied original excerpt.} The original text begins with the fragments ``Is $G$ cocompactly cubulable?'', ``Is $G$ necessarily residually finite?'', and a question about removing ``uniform'' from an unavailable preceding statement. It then contains a different passage headed ``Problem 10. Covers of cube complexes, Shepherd.'' The supplied curation instead associates the main questions with Problems 11 and 12. The missing antecedents and this extraction mismatch prevent verification that (A)--(C) reproduce the original questions exactly.

For preservation of the additional mathematical content, the Problem 10 passage asks:
\begin{itemize}
\item[(i)] Attributed to Haglund: if $X_1$ and $X_2$ are finite special cube complexes with a common universal cover, must they have a common finite cover?
\item[(ii)] Is there an algorithm which, given two finite nonpositively curved cube complexes, determines whether their universal covers are isomorphic?
\end{itemize}
Here a special cube complex is understood in the usual sense of a nonpositively curved cube complex admitting a cubical local isometry to a Salvetti complex. A common finite cover means a cube complex admitting finite-sheeted covering maps to both complexes. The precise category of universal-cover isomorphisms in (ii) is not specified in the excerpt.

The excerpt reports affirmative answers to (i) for graphs, attributed to Leighton's theorem; for complexes with free fundamental groups, attributed to Bridson--Shepherd; and when the common universal cover is the right-angled building associated to a graph product of finite groups, attributed to Shepherd. It also states that the specialness assumption is necessary, citing counterexamples without it whose universal cover is a product of trees, using examples of Wise or Burger--Mozes. For (ii), it reports an algorithm for graphs, attributed to Angluin and based on colour refinement. These are supplied contextual claims, not independently verified results here. Their presence does not establish that Problem 10 is an intended additional target of permanent ID 30006622.
\subsection{Short English statement}
The curated record asks whether elementwise cubical actions force a hyperbolic group to admit a proper cocompact cubical action or a proper affine Hilbert-space action; whether proper cocompact actions on finite-rank median spaces force residual finiteness or linearity; and whether local medianity implies global medianity without a uniform local radius. Several conventions and missing hypotheses remain unresolved. The original excerpt also includes separate cube-complex covering questions, whose relationship to this record is unclear. This is a statement-only draft requiring human mathematical review.
\subsection{Sources}
\begin{itemize}
\item Supplied local UnsolvedMath registry record, permanent ID \texttt{30006622}, exact source problem code \texttt{catalogue entry}, titled ``Hyperbolic Cubulability, Residual Finiteness, and Local Medianity,'' in the collection ``Oberwolfach Reports Open Problems.'' No proposer is identified for the main bundled questions.
\item Supplied source citation: \emph{Median Geometry and Applications} (2026), \emph{Oberwolfach Reports}, pp.~487--540, DOI \texttt{10.4171/owr/2026/8}. Supplied source link: \url{https://doi.org/10.4171/owr/2026/8}. The background describes Oberwolfach Reports as documenting mathematical workshops and their open-problem discussions.
\item Supplied literature source labelled ``OWR report'': \url{https://ems.press/content/serial-article-files/53603}.
\item Supplied literature source labelled ``Bowditch preprint'': \url{https://bhbowditch.com/papers/ch-median.pdf}. No title, version, or precise supporting passage is supplied.
\item The supplied literature assessment was checked on exactly \texttt{2026-08-22}. It labels the record open and states that the extraction bundles unresolved cubulability, residual-finiteness, and non-uniform local-median questions from Problems 11 and 12, with no complete resolution found. This is supplied dated curation, described in the record as literature triage derived from the OWR curation; it has not been independently verified here. The supplied record expressly calls for independent verification of time-sensitive status and cited claims.
\item The original excerpt attributes its Problem 10 heading to Shepherd, its first covering question to Haglund, and contextual results or examples to Leighton, Bridson--Shepherd, Shepherd, Wise, Burger--Mozes, and Angluin. No separate bibliographic citations for these attributions are supplied.
\item The registry carries the supplied label \texttt{corrected\_verified} and says that only formula delimiters, notation, grammar, and typography were repaired while mathematical content and scope were preserved. This label is a source assertion and does not resolve the mismatch between the original excerpt and the curated statement.
\end{itemize}
\end{document}
